\documentclass[11pt]{article}
\usepackage[T1]{fontenc}
\usepackage[margin=1in]{geometry}
\usepackage{amsmath,amssymb,mathtools}
\usepackage[colorlinks=true,linkcolor=blue,citecolor=blue,urlcolor=blue]{hyperref}

\title{Perturbative action around AdS$_3$}
\author{}
\date{}

\newcommand{\dd}{\mathrm{d}}
\newcommand{\cO}{\mathcal{O}}
\newcommand{\Lg}{\mathcal{L}_{\mathrm g}}

\begin{document}
\maketitle

This fragment fixes the perturbative action used in the Noether-charge
calculation. We use the normalization
\begin{align}
  \lambda^2=8\pi G,
\end{align}
so that the gravitational part of the action is
\begin{align}
  I_{\mathrm{grav}}[g]
  =\frac{1}{2\lambda^2}\int_M \dd^3x\,\sqrt{-g}\,(R+2).
\end{align}
For the scalar example,
\begin{align}
  I[g,\phi]
  =\frac{1}{2\lambda^2}\int_M \dd^3x\,\sqrt{-g}\,(R+2)
  -\frac12\int_M \dd^3x\,\sqrt{-g}\,
  \left(\nabla_\mu\phi\nabla^\mu\phi+m^2\phi^2\right).
\end{align}
The AdS radius is set to one. The background metric $g^{(0)}_{\mu\nu}$
satisfies
\begin{align}
  R^{(0)}_{\mu\nu}=-2g^{(0)}_{\mu\nu},
  \qquad
  R^{(0)}=-6.
\end{align}
All indices below are raised and lowered with $g^{(0)}_{\mu\nu}$, and
$\nabla^{(0)}_\mu$ denotes the corresponding background covariant derivative.

We expand
\begin{align}
  g_{\mu\nu}
  &=g^{(0)}_{\mu\nu}+\lambda h_{\mu\nu}
    +\lambda^2 k_{\mu\nu}+\cO(\lambda^3),\\
  \phi&=\varphi+\cO(\lambda).
\end{align}
The traces are
\begin{align}
  h=g^{(0)\mu\nu}h_{\mu\nu},
  \qquad
  k=g^{(0)\mu\nu}k_{\mu\nu}.
\end{align}
For any symmetric tensor $X_{\mu\nu}$, write
$X=g^{(0)\mu\nu}X_{\mu\nu}$. The inverse metric and volume density are
\begin{align}
  g^{\mu\nu}
  &=g^{(0)\mu\nu}
    -\lambda h^{\mu\nu}
    +\lambda^2\left(h^\mu{}_\rho h^{\rho\nu}-k^{\mu\nu}\right)
    +\cO(\lambda^3),\\
  \sqrt{-g}
  &=\sqrt{-g^{(0)}}\left[
      1+\frac{\lambda}{2}h
      +\lambda^2\left(
        \frac12 k-\frac14 h_{\mu\nu}h^{\mu\nu}+\frac18 h^2
      \right)
    \right]+\cO(\lambda^3).
\end{align}

At the level of the action, after discarding total-derivative terms in the
quadratic coefficient, write
\begin{align}
  \sqrt{-g}(R+2)
  =\sqrt{-g^{(0)}}\left[
    \Lg^{(0)}
    +\lambda \Lg^{(1)}[h]
    +\lambda^2\left(\Lg^{(1)}[k]+\Lg^{(2)}[h]\right)
  \right]+\cO(\lambda^3),
\end{align}
where
\begin{align}
  \Lg^{(0)}&=-4,\\
  \Lg^{(1)}[X]
  &=\nabla^{(0)}_\mu\nabla^{(0)}_\nu X^{\mu\nu}
    -\nabla^{(0)2}X,
\end{align}
and
\begin{align}
  \Lg^{(2)}[X]
  &=\frac12 X^2-X_{\mu\nu}X^{\mu\nu}
    +X^{\mu\nu}\nabla^{(0)}_\mu\nabla^{(0)}_\nu X
    -\frac14\nabla^{(0)}_\mu X\nabla^{(0)\mu}X \notag\\
  &\quad
    -\nabla^{(0)}_\mu X^{\mu\nu}
      \nabla^{(0)}_\rho X_\nu{}^\rho
    +\nabla^{(0)\mu}X
      \nabla^{(0)}_\nu X_\mu{}^\nu
    -X^{\mu\nu}
      \nabla^{(0)}_\nu\nabla^{(0)}_\rho X_\mu{}^\rho \notag\\
  &\quad
    -X^{\mu\nu}
      \nabla^{(0)}_\rho\nabla^{(0)}_\nu X_\mu{}^\rho
    +\frac12 X
      \nabla^{(0)}_\mu\nabla^{(0)}_\nu X^{\mu\nu}
    +X^{\mu\nu}\nabla^{(0)2}X_{\mu\nu} \notag\\
  &\quad
    -\frac12 X\nabla^{(0)2}X
    -\frac12\nabla^{(0)}_\mu X_{\nu\rho}
      \nabla^{(0)\rho}X^{\mu\nu}
    +\frac34\nabla^{(0)}_\rho X_{\mu\nu}
      \nabla^{(0)\rho}X^{\mu\nu}.
\end{align}

Therefore the full action has the expansion
\begin{align}
  I[g,\phi]
  &=\frac{1}{2\lambda^2}
    \int_M\dd^3x\,\sqrt{-g^{(0)}}\,\Lg^{(0)}
    +\frac{1}{2\lambda}
    \int_M\dd^3x\,\sqrt{-g^{(0)}}\,\Lg^{(1)}[h]
    \notag\\
  &\quad
    +I^{(0)}[h,k,\varphi]+\cO(\lambda),
\end{align}
with finite part
\begin{align}
  I^{(0)}[h,k,\varphi]
  &=\frac12\int_M\dd^3x\,\sqrt{-g^{(0)}}
    \left(\Lg^{(1)}[k]+\Lg^{(2)}[h]\right)
    \notag\\
  &\quad
    -\frac12\int_M\dd^3x\,\sqrt{-g^{(0)}}
    \left(\nabla_\mu\varphi\nabla^\mu\varphi+m^2\varphi^2\right).
\end{align}

The Noether-charge calculation in the next fragment uses only the finite
action $I^{(0)}[h,k,\varphi]$. The divergent, order-$\lambda^{-1}$, and
discarded total-derivative pieces are background or boundary contributions in
this perturbative expansion.

\end{document}
